\scanpage{83}
% archive sheet number: 82
% Long oblique strokes cancel this draft.
\deleted{IV 7.5.}

\deleted{Je ne suis pas très content encore de la dernière version, car on n'énonce pas ce qu'on démontre vraiment, et par conséquent il y a des cas (tels la ``$\operatorname{copfof}\quad n$'', ou ``intègre'') qui échappent à l'énoncé tel qu'il est donné, mais sont inclus dans la démonstration. Aussi je propose de réénoncer les choses ainsi:}

\deleted{\uline{Théorème} 7.5.1.\quad Soit $A\quad B$ un homomorphisme local d'anneaux locaux noéthériens, $\text{\uuline{R}}(C)$ une propriété \deleted{d'anneaux locaux} faisant intervenir un argument qui est un anneau local \unsure{noehhérien|noethérien}. On suppose\deleted{\unsure{r|x}} que $(B,R)$ satisfait à :}

\deleted{a) Pour tout anneau \inserted{local} $B^{\deleted{\unsure{'|x}}}$ localise d'une $B$-algèbre finie, et tout élément régulier $t\quad\underline{r}(C)$, $\text{\uuline{R}}(C/tC)$ implique $\text{\uuline{R}}(C)$.}

\deleted{b) Si $C$ est une $B$-algèbre finie, alors $U_{\text{\uuline{R}}}(\operatorname{Spec}(C))$ est ouvert.}

\deleted{On suppose de plus que $A$ satisfait les conditions suivantes\deleted{\unsure{x}} (remplies si $A$ est ``excellent'' (7.8.2.) en particulier si $A$ est complet):}

\deleted{\handwritten{$\alpha$}) $A$ est universellement japonais.}

\deleted{\handwritten{$\beta$}) Pour toute algèbre finie intègre $C$ sur $A$, $\operatorname{Spec}(C)$ contient un ouvert non vide qui est régulier,}

\deleted{Soit $k$ le corps résiduel de $A$, \inserted{et supposons le corps résiduel de $B$ fini sur celui de $A$.} Supposons $B$ plat sur $A$ et que $B\otimes_A k$ sur $k$ satisfait la condition ``$Q$ géométrique''; Alors toutes les fibres de $\operatorname{Spec}(B)$ sur $\operatorname{Spec}(A)$ satisfont la condition $Q$ géométrique.}

\deleted{Démonstration. Les hypothèses faites sur $A\quad B$ sont stables par changement de base fini $A\quad A'$, pourvu qu'on prenne soin de reformuler les hypothèses pour $A,B$ semi-locaux, l'homomorphisme $A\quad B$ semi-local (i.e. l'image inverse d'un idéal maximal est un idéal maximal) et qu'on remplace le corps résiduel de $A$ par les corps résiduels de $A$. \deleted{(NB alors} (NB cette extension ne mérite pas de figurer dans l'énoncé du théorème, c'est juste un intermédiaire technique pour la démonstration). On est ainsi ramené au cas $A$ intègre et à voir que la fibre générique satisfait}
